\section{Conclusion}\label{sec:conclusion}

A decade of empirical work has established that restricting a procurement auction to small firms raises winning prices. This paper goes a step further: it identifies which mechanism produces the markup under explicit modeling assumptions and prices the welfare consequence against a specific policy alternative. The March 2018 extension of S\~ao Paulo's SME-only regime to medical supplies supplies the identifying variation; Preg\~ao drop-out bids identify the type-specific cost distributions under \citet{haile2003}; auction-level common scale is absorbed by Krasnokutskaya-style shrinkage; entry is handled in the \citet{atheyseira2011} spirit. The exercise runs under observed equilibrium entry, with strict primitive invariance reported as a bound rather than the main specification.

Five points organize the contribution. First, the difference-in-differences benchmark of Appendix~\ref{app:did} establishes a 10--11 percent rise in absolute winning prices, the average winning-price response that anchors the structural calibration. The structural counterfactual, on the structural normalization $p_{S_3} - p_{S_1}$ scaled by $p_{S_1}$, is roughly three to four times that DiD coefficient (about 34 percent in non-pharma and 47 percent in pharma); the implied total welfare cost (allocative DWL plus MCPF distortion at $\lambda = 0.30$) reaches 28.7 percent of $p_{S_1}$ in non-pharma and 47.0 percent in pharma. Second, the simulated set-aside price effect operates predominantly within the auction under the model's main specification: the within-auction component---which I label \emph{sheltered bidding}---accounts for roughly three-quarters of the simulated total effect, in a window stable across the cost-distribution estimators reported. Third, endogenous SME entry is the government's partial offset rather than the source of the markup: a fixed-pool counterfactual would deliver a simulated price effect roughly 50--60 percent larger than the realized one. Fourth, a 10 percent price preference welfare-dominates the set-aside in thick, standardized non-pharmaceutical markets across both the main and the strict-invariance specifications, at near-zero fiscal cost. Fifth, in thin, heterogeneous pharmaceutical markets the ranking depends on how the post-policy SME pool is modeled; the conditional reading is itself a substantive result about where bidder exclusion delivers proportional returns and where it does not.

The policy reading follows. The question is not whether bidder exclusion is universally good or bad, but when: when the protected pool is thick enough, the good standardized enough, and the planner's welfare weight on SME surplus low enough that a moderate price preference recovers nearly the same distributive gain at a fraction of the welfare cost. Where the pool is thin and the good heterogeneous, the choice depends on equilibrium-selection details this paper identifies but does not resolve. A general defense of bidder-exclusion rules across procurement categories is harder to sustain than current Brazilian and international practice would suggest. A uniform replacement of exclusion by preference is stronger than the data support. The choice belongs at the level of the product class, not the level of the procurement code.
